\documentclass[11pt]{article}

\usepackage[margin=1in]{geometry}
\usepackage{amsmath,amssymb,amsfonts,bm}
\usepackage{graphicx}
\usepackage{xcolor}
\usepackage{float}
\usepackage{enumitem}
\usepackage{booktabs}
\usepackage{hyperref}
\usepackage{xurl}
\usepackage{fancyhdr}
\usepackage{xcolor}
\usepackage{listings}

\lstset{
	basicstyle=\ttfamily\small,
	breaklines=true,
	columns=fullflexible
}
% ---------- URL and Hyperlink setup ----------
\Urlmuskip=0mu plus 1mu
\hypersetup{
	breaklinks=true,
	colorlinks=true,
	linkcolor=blue,
	urlcolor=blue
}

% ---------- Header setup ----------
\pagestyle{fancy}
\fancyhf{}
\fancyhead[L]{\today}
\fancyhead[C]{ECE 506: Homework \#6}
\fancyhead[R]{Scott Nguyen}
\renewcommand{\headrulewidth}{0.4pt}
\fancyfoot[C]{\thepage}

% ---------- Remove paragraph indent ----------
\setlength{\parindent}{0pt}

\begin{document}
	
	\noindent\textbf{Problem 1. Basic Concepts}
	
	\vspace{0.5em}
	For this problem, the goal is to explore basic concepts using the Neural Network
	playground: \url{https://playground.tensorflow.org/}.
	
	\vspace{1em}
	\noindent\textbf{1(a).}
	
	\begin{enumerate}[label=1.a.\roman*., leftmargin=1.2cm, itemsep=0.4em]
		\item What do the orange and blue colors actually represent for the data input, the features, the hidden layers, and the output?
		
		\noindent\textbf{Solution} 
		\begin{itemize}
			\item \textbf{Data input:} Orange = class 1, Blue = class 0.
			\item \textbf{Features:} Orange = feature $> 0$, Blue = feature $< 0$, White = $\approx 0$.
			\item \textbf{Hidden layers:} Orange = positive activation, Blue = negative activation, White = activation $\approx 0$.
			\item \textbf{Output:} Orange = predicted class 1, Blue = predicted class 0.
		\end{itemize}
		
		
		\item Change the activation function to ReLU. What does the white color represent for the data input, the features, the hidden layers, and the output?
		
		\noindent\textbf{Solution} 
		\begin{itemize}
			\item \textbf{White} indicates regions where the neuron output is exactly zero: $\text{ReLU}(z) = 0$.
		\end{itemize}
		
		\item Write down the equation of the Neural Network represented by the default values.
		
		\noindent\textbf{Solution}
		Let $\mathbf{x} = [x,\, y]^T$. The network has two hidden layers with $\tanh$ activation.
		\[
		h^{(1)} = \tanh\!\left(W^{(1)} \mathbf{x} + b^{(1)}\right)
		\]
		\[
		h^{(2)} = \tanh\!\left(W^{(2)} h^{(1)} + b^{(2)}\right)
		\]
		\[
		\hat{y} = \sigma\!\left(W^{(3)} h^{(2)} + b^{(3)}\right)
		\]
		\[
		\text{where } \sigma(z) = \frac{1}{1 + e^{-z}}.
		\]
		
	\end{enumerate}
	
	\vspace{1em}
	\noindent\textbf{1(b).} Train the Neural Network and answer the following:
	
	\begin{enumerate}[label=1.b.\roman*., leftmargin=1.2cm, itemsep=0.4em]
		\item What does it mean to train through an epoch?
		
		\noindent\textbf{Solution}
		Training through one epoch means using every training sample exactly once to update the network parameters.
		
		\item How do you know that the network converged in terms of the test loss? Provide screenshots.
		
		\noindent\textbf{Solution}
		Yes. The network has converged because the test loss stabilizes at a small value
		(\texttt{0.003}) and no longer decreases over additional epochs. The decision
		boundary is also stable. The screenshot in Fig.~\ref{fig:converged} shows the
		final test loss and prediction region.
		
		\begin{figure}[H]
			\centering
			\includegraphics[width=0.4\textwidth]{plots/fig_1.png}
			\caption{Test loss and decision boundary at convergence.}
			\label{fig:converged}
		\end{figure}
		
		\item What is the minimum number of epochs that it takes to converge? Explain your answer and provide screenshots.
		
		\noindent\textbf{Solution}
		The minimum number of epochs required for convergence is the point after which
		the test loss no longer decreases. As shown in Fig.~\ref{fig:epoch_converged},
		the test loss stabilizes around epoch $\approx 170$, indicating convergence.
		
		\begin{figure}[H]
			\centering
			\includegraphics[width=0.7\textwidth]{plots/fig_2.png}
			\caption{Test loss curve and decision boundary near convergence (epoch $\approx 173$).}
			\label{fig:epoch_converged}
		\end{figure}
		
		\item Increase the learning rate by a factor of 100. What happened to the convergence? Explain. Provide screenshots and discuss convergence in terms of the number of epochs.
		
		\noindent\textbf{Solution}
		Increasing the learning rate by a factor of 100 makes training unstable. As shown
		in Fig.~\ref{fig:high_lr}, the test loss oscillates instead of decreasing smoothly,
		and the decision boundary becomes irregular. The model does not converge within
		a reasonable number of epochs.
		
		\begin{figure}[H]
			\centering
			\includegraphics[width=0.7\textwidth]{plots/fig_3.png}
			\caption{Training with learning rate increased by a factor of 100. The test loss oscillates and does not converge.}
			\label{fig:high_lr}
		\end{figure}
		
		\item Repeat the previous question by reducing the learning rate to the lowest value.
		
		\noindent\textbf{Solution}
		With the learning rate reduced to the lowest value, convergence becomes extremely
		slow. As shown in Fig.~\ref{fig:low_lr}, both the test loss and decision boundary
		improve only gradually, requiring many more epochs. The model eventually moves in
		the correct direction, but at a very slow rate.
		
		\begin{figure}[H]
			\centering
			\includegraphics[width=\textwidth]{plots/fig_4.png}
			\caption{Training with the lowest learning rate. Convergence is very slow and the test loss remains high for many epochs.}
			\label{fig:low_lr}
		\end{figure}
		
	\end{enumerate}
	
	\newpage
	
	\noindent\textbf{Problem 2. Backpropagation}
	
	\vspace{0.5em}
	\noindent\textbf{2(a).} For this problem, we will follow the tutorial on autograd given at: \href{https://docs.pytorch.org/tutorials/beginner/blitz/autograd_tutorial.html#sphx-glr-beginner-blitz-autograd-tutorial-py}{PyTorch autograd}.\\
	  
	We use a loss function to train neural networks using a $l(.)$ that takes a vector and outputs a scalar as in traditional optimization methods of studying $f(x)$. However, neural networks use function composition to propagate the inputs tothe outputs. In the autograd example, the notation is as follows:
	
	\begin{itemize}
		\item \(l(\cdot)\):  is the output scalar loss function (e.g., mean squared error).
		\item \(y\): denotes the top output layers of the neural network.
		\item \(x\): denotes the input to the neural network.
	\end{itemize}
	
	In order to train Neural-networks using the loss function, we need to compute gradients of $l(.)$ with respect to all of the variables. In the simplest form, we use functional composition to express $l(.)$ for a specific input: $l(y(x))$. When designing neural-networks, we can assume that the derivatives of each level are available to us. This means that we assume that the following derivatives are given to us:
	
	\[
	\frac{\partial l}{\partial y_i}, \qquad \frac{\partial y_i}{\partial x_j}
	\quad \text{for } i = 1,\dots,m \text{ and } j = 1,\dots,n.
	\]
	
	We then have to determine all other derivatives. To determine the dependence
	of the loss-function based on the input, we use:
	
	\[
	\frac{\partial l}{\partial x} 
	= 
	\begin{bmatrix}
		\frac{\partial y_1}{\partial x_1} & \cdots & \frac{\partial y_m}{\partial x_1} \\
		\vdots & \ddots & \vdots \\
		\frac{\partial y_1}{\partial x_n} & \cdots & \frac{\partial y_m}{\partial x_n}
	\end{bmatrix}
	\begin{bmatrix}
		\frac{\partial l}{\partial y_1} \\ \vdots \\ \frac{\partial l}{\partial y_m}
	\end{bmatrix}
	\]
	
	Multiply out the matrix by the vector to derive an expression for $\frac{\partial l}{\partial x_i}$. How many terms do you have?
	
	\noindent\textbf{Solution}
	The $i$-th component of $\dfrac{\partial l}{\partial x}$ is obtained by taking the
	dot product of the $i$-th row of the Jacobian $\dfrac{\partial y}{\partial x}$ with
	the vector $\dfrac{\partial l}{\partial y}$:
	\[
	\frac{\partial l}{\partial x_i}
	= \sum_{k=1}^{m}
	\frac{\partial y_k}{\partial x_i}\,
	\frac{\partial l}{\partial y_k},
	\qquad i = 1,\dots,n.
	\]
	There are $m$ total terms in the expression for $\frac{\partial l}{\partial x_i}$, one for each output component $y_k$.
	
	
	\vspace{0.5em}
	\noindent\textbf{2(b).} Let us introduce one more layer. Suppose that the loss is given by \(l(y(g(x)))\). Apply 2(a) recursively to derive an expression for \(\partial l / \partial x_i\) for this case. Assume that \(g\) outputs a \(p\)-dimensional vector that serves as the input to \(y\).
	Assume that all other dimensions remain as given in 1(a). You will need to clearly label all matrix dimensions in your derivation.
	
	\noindent\textbf{Solution}
	Let $x \in \mathbb{R}^n$, $g(x) \in \mathbb{R}^p$, and $y(g) \in \mathbb{R}^m$.
	Then
	\[
	\frac{\partial l}{\partial y} \in \mathbb{R}^{1 \times m}, \quad
	\frac{\partial y}{\partial g} \in \mathbb{R}^{m \times p}, \quad
	\frac{\partial g}{\partial x} \in \mathbb{R}^{p \times n}.
	\]
	By the chain rule,
	\[
	\frac{\partial l}{\partial g}
	= \frac{\partial l}{\partial y} \frac{\partial y}{\partial g}
	\in \mathbb{R}^{1 \times p},
	\qquad
	\frac{\partial l}{\partial x}
	= \frac{\partial l}{\partial g} \frac{\partial g}{\partial x}
	\in \mathbb{R}^{1 \times n}.
	\]
	Component-wise,
	\[
	\frac{\partial l}{\partial x_i}
	= \sum_{j=1}^{p} \sum_{k=1}^{m}
	\frac{\partial l}{\partial y_k}
	\frac{\partial y_k}{\partial g_j}
	\frac{\partial g_j}{\partial x_i}.
	\]
	
	\noindent\textbf{2(c).} Repeat 1(b) for \(l\big(y(g_{2}(g_{1}(x)))\big)\). Assume that the
	output of \(g_{2}\) is \(p_{2}\)-dimensional and the output of \(g_{1}\) is \(p_{1}\)-dimensional.
	Make sure to clearly indicate all matrix dimensions.
	\noindent\textbf{Solution}
	Let $x \in \mathbb{R}^n$, $g_{1}(x) \in \mathbb{R}^{p_{1}}$, $g_{2}(g_{1}) \in \mathbb{R}^{p_{2}}$,
	and $y(g_{2}) \in \mathbb{R}^m$, with scalar loss $l(y)$.
	
	Then the relevant Jacobians are:
	\[
	\frac{\partial l}{\partial y} \in \mathbb{R}^{1 \times m}, \quad
	\frac{\partial y}{\partial g_{2}} \in \mathbb{R}^{m \times p_{2}}, \quad
	\frac{\partial g_{2}}{\partial g_{1}} \in \mathbb{R}^{p_{2} \times p_{1}}, \quad
	\frac{\partial g_{1}}{\partial x} \in \mathbb{R}^{p_{1} \times n}.
	\]
	
	By the chain rule,
	\[
	\frac{\partial l}{\partial x}
	= \frac{\partial l}{\partial y}
	\frac{\partial y}{\partial g_{2}}
	\frac{\partial g_{2}}{\partial g_{1}}
	\frac{\partial g_{1}}{\partial x}
	\in \mathbb{R}^{1 \times n}.
	\]
	
	Component-wise, for $i = 1,\dots,n$,
	\[
	\frac{\partial l}{\partial x_i}
	= \sum_{k=1}^{m} \sum_{j=1}^{p_{2}} \sum_{r=1}^{p_{1}}
	\frac{\partial l}{\partial y_k}
	\frac{\partial y_k}{\partial (g_{2})_j}
	\frac{\partial (g_{2})_j}{\partial (g_{1})_r}
	\frac{\partial (g_{1})_r}{\partial x_i}.
	\]
	\newpage
	
	\vspace{0.5em}
	\noindent\textbf{Problem 3. Automatic Differentiation}
	
	The tutorial on autograd in \texttt{autograd\_tutorial.ipynb} shows an example of how to
	verify that automatic differentiation works for a simple function. Modify the
	example to compute
	
	\[
	\frac{\partial}{\partial x_i}
	\left[
	\frac{1}{6}
	\sum_{i=1}^{6} (x_i - 2)^2 + x_i
	\right].
	\]
	
	Verify that your example works by computing the tensor output for
	\(\partial / \partial x_i\) when \(x_i = 1\).
	
	\noindent\textbf{Solution.}
	We consider the function
	\[
	f(x) = \frac{1}{6}\sum_{i=1}^{6}(x_i - 2)^2 + x_i.
	\]
	The derivative of the $i$th term is
	\[
	\frac{\partial f}{\partial x_i}
	= \frac{1}{3}(x_i - 2) + 1.
	\]
	Evaluating at $x_i = 1$ gives
	\[
	\frac{\partial f}{\partial x_i}\bigg|_{x_i=1} = \frac{2}{3}.
	\]
	
	To verify this using \texttt{torch.autograd}, we define the tensor 
	\texttt{x = torch.ones(6, requires\_grad=True)} and compute:
	\[
	\texttt{f = (1/6) * torch.sum((x - 2)**2) + torch.sum(x)}, \qquad
	\texttt{f.backward()}.
	\]
	The resulting gradient is
	\[
	\nabla f(x) = 
	\begin{bmatrix}
		0.6667 & 0.6667 & 0.6667 & 0.6667 & 0.6667 & 0.6667
	\end{bmatrix},
	\]
	which matches the analytic value $\frac{2}{3}$ for each component.
	
	\newpage
	
	\vspace{0.5em}
	\noindent\textbf{Problem 4. The Basic Neural Network Interface}
	
	For this problem, please study \texttt{fahrenheit\_to\_celsius\_corrected.ipynb} and\\ \texttt{neural\_networks\_tutorial.ipynb}. The SGD implementation is given in:\\
	\url{https://pytorch.org/docs/stable/optim.html}
	
	\vspace{0.5em}
	\noindent\textbf{4(a).} Modify the example and show how to compute the conversion
	from Fahrenheit to Celsius. This is the inverse of the conversion given in the example.
	
	\noindent\textbf{Solution.}
	To learn the inverse temperature conversion, we fit the linear model
	\[
	C = \alpha_0 + \alpha_1 F
	\]
	using autograd and gradient descent. Random Fahrenheit samples are generated,
	their true Celsius values are computed from
	\[
	C_{\text{true}} = \frac{F - 32}{1.8},
	\]
	and the parameters \(\alpha_0,\alpha_1\) are optimized by minimizing the 
	mean–squared error. The training converges to the exact inverse mapping.
	
	\begin{verbatim}
		Fahrenheit to Celsius
		True:    C = -17.78 + 0.56·F
		Learned: C = -17.78 + 0.56·F
		Final loss: 4.411049303598702e-11
		Iterations: 2000
	\end{verbatim}
	
	\vspace{0.5em}
	\noindent\textbf{4(b).} Translate the example into the neural-network framework and
	learn the parameters using:
	
	\begin{itemize}
		\item Stochastic Gradient Descent (SGD) without momentum,
		\item Stochastic Gradient Descent (SGD) with momentum.
	\end{itemize}
	
	Based on the documentation in \url{https://pytorch.org/docs/stable/optim.html}, 
	the basic SGD update is
	
	\[
	p_{t+1} = p_{t} - \text{lr} \cdot \nabla_p f,
	\]
	
	where \(\text{lr}\) is the learning rate and \(p\) is a parameter being optimized.
	
	With momentum, for step \(t+1\), the approach is modified to:
	
	\[
	v_{t+1} = \mu \cdot v_t + g_{t+1},
	\qquad
	p_{t+1} = p_t - \text{lr} \cdot v_{t+1},
	\]
	
	where \(g_{t+1} = \nabla_p f_{t+1}\).  
	You do not need to implement SGD manually—simply call the appropriate optimizer
	functions in PyTorch. Document your results for several different parameter choices.
	
	\noindent\textbf{Solution.}
	We translate the Fahrenheit-to-Celsius conversion into the neural-network
	framework by using a single fully connected layer \(\texttt{nn.Linear(1,1)}\),
	which implements the model
	\[
	C = \alpha_0 + \alpha_1 F.
	\]
	Using the same dataset as in part 4(a), we train this network twice with
	PyTorch's \texttt{optim.SGD}: first with standard SGD (no momentum), then with
	SGD including momentum (\(\mu = 0.9\)). In both cases we minimize the
	mean–squared error between predicted and true Celsius values.
	
	For both optimizers, the learned parameters converge to
	\[
	C \approx -17.78 + 0.56\,F,
	\]
	matching the true inverse conversion. The final printed results are:
	
	\begin{verbatim}
		Problem 4(b): Neural network with SGD
		SGD (no momentum)
		Learned: C = -17.78 + 0.56·F
		Final loss: 4.5929482439532876e-11
		
		SGD (with momentum = 0.9)
		Learned: C = -17.78 + 0.56·F
		Final loss: 8.45830089996058e-12
	\end{verbatim}
	
	\newpage
	
	\vspace{0.5em}
	\noindent\textbf{Problem 5. Adam Optimization}
	
	This problem is based on \textit{ADAM: A Method for Stochastic Optimization}, which can be accessed at
	\url{https://arxiv.org/pdf/1412.6980}. Additionally, consult the \texttt{ML-intro.pdf} document.
	
	\vspace{0.5em}
	\noindent\textbf{5(a).} Consider the update for the mean that is given by the Adam algorithm.
	Use the definition of the sample mean to derive an expression for \(\beta_{1}\) that makes
	the bias update correct.
	
	\noindent\textbf{Solution}\\
	Adam updates the first moment using
	\[
	m_t = \beta_1 m_{t-1} + (1-\beta_1) g_t.
	\]
	The unbiased sample mean of the gradients is
	\[
	\bar{g}_t = \frac{1}{t}\sum_{i=1}^t g_i.
	\]
	Expanding the Adam update gives
	\[
	m_t = (1-\beta_1)\sum_{i=1}^t \beta_1^{t-i} g_i.
	\]
	To match the uniform weight \(1/t\) in the sample mean, we require
	\[
	1-\beta_1 = \frac{1}{t},
	\]
	so the correct choice is
	\[
	\beta_1 = 1 - \frac{1}{t}.
	\]
	
	\vspace{0.5em}
	\noindent\textbf{5(b).} Does your expression for \(\beta_{1}\) in 5(a) depend on the batch size \(n\)?
	What happens when you double the batch size?
	
	\noindent\textbf{Solution}\\
	The expression derived in part 5(a),
	\[
	\beta_1 = 1 - \frac{1}{t},
	\]
	depends only on the iteration index \(t\) and does not involve the batch size \(n\).
	Doubling the batch size changes the variance of each gradient estimate but does not
	change the form of the bias in the exponential moving average. Thus, the value of
	\(\beta_1\) required to correct the bias remains the same regardless of batch size.
	
	\vspace{0.5em}
	\noindent\textbf{5(c).} Consider the update for the sample variance given by the Adam algorithm.
	
	\begin{enumerate}[label=5.c.\roman*., leftmargin=1.4cm, itemsep=0.4em]
		\item Use the definition of the sample variance to derive an expression for \(\beta_{2}\).\\
		\noindent\textbf{Solution}\\
		Adam updates the second moment via
		\[
		v_t = \beta_2 v_{t-1} + (1-\beta_2) g_t^2.
		\]
		The unbiased sample variance estimate of the squared gradients is
		\[
		\frac{1}{t}\sum_{i=1}^t g_i^2.
		\]
		Expanding the Adam update gives
		\[
		v_t = (1-\beta_2)\sum_{i=1}^t \beta_2^{\,t-i} g_i^2.
		\]
		To match the uniform weight \(1/t\) used in the sample variance, we require
		\[
		1-\beta_2 = \frac{1}{t},
		\]
		which yields
		\[
		\beta_2 = 1 - \frac{1}{t}.
		\]
		Thus, the same bias condition applies to the second-moment estimate.
		
		\item What did you assume for the sample mean in order to derive your expression?\\
		\noindent\textbf{Solution}\\
		The derivation assumed that the sample mean used in the variance formula was the
		unbiased average
		\[
		\bar{g}_t = \frac{1}{t}\sum_{i=1}^t g_i,
		\]
		so that both the mean and variance rely on uniform weights \(1/t\).
		
		\item Why do we need a bias correction term for the second moment?\\
		\noindent\textbf{Solution}\\
		Without bias correction, the exponential moving average
		\(
		v_t = \beta_2 v_{t-1} + (1-\beta_2) g_t^2
		\)
		is biased low during early iterations because it begins at zero and places most
		weight on that initial value. The bias correction term removes this initialization
		bias so that \(v_t\) becomes an unbiased estimate of the true second moment.
		
		\item Is the assumption in 5.c.i correct at the initial epochs?\\
		\noindent\textbf{Solution}\\
		No. In the early epochs only a few gradients are available, so the sample mean is
		not yet accurate. This makes the assumption in 5.c.i invalid at the start.
		
		\item How about when you approach convergence?\\
		\noindent\textbf{Solution}\\
		Near convergence, many gradients have been observed, so the sample mean becomes
		accurate and the assumption in 5.c.i holds much more closely.
		
	\end{enumerate}
	
	\vspace{0.5em}
	\noindent\textbf{5(d).} What happens to the variance of the sample mean if we double the batch
	size?\\
	\noindent\textbf{Solution}\\
	Doubling the batch size reduces the variance of the sample mean by a factor of two,
	since averaging over twice as many samples cuts the variance in half.
	
	\newpage
	
	\vspace{0.5em}
	\noindent\textbf{5(e).} Give the general shape of the \(n\)-dimensional covariance matrix that the
	Adam algorithm assumes. When is this assumption violated?\\
	\noindent\textbf{Solution}\\
	Adam assumes that the covariance matrix of the gradients is diagonal, i.e.,
	\[
	\Sigma = 
	\begin{bmatrix}
		\sigma_1^2 & 0 & \cdots & 0 \\
		0 & \sigma_2^2 & \cdots & 0 \\
		\vdots & \vdots & \ddots & \vdots \\
		0 & 0 & \cdots & \sigma_n^2
	\end{bmatrix}.
	\]
	This assumes all gradient components are uncorrelated. The assumption is violated
	when parameters exhibit significant cross-correlations, causing nonzero
	off-diagonal terms in the true covariance matrix.
	
	\vspace{0.5em}
	\noindent\textbf{5(f).} Suppose that the variance of the gradient for a neural network weight
	component is very small. Derive a simplified expression for the final expression of parameter update.\\
	\noindent\textbf{Solution}\\
	In Adam, the parameter update is
	\[
	\theta_{t+1}
	= \theta_t
	- \alpha \frac{\hat{m}_t}{\sqrt{\hat{v}_t} + \varepsilon}.
	\]
	If the gradient variance is very small, then
	\(
	v_t \approx g_t^2
	\)
	and
	\(
	\sqrt{\hat{v}_t} \approx |g_t|.
	\)
	Thus the update becomes
	\[
	\theta_{t+1}
	\approx \theta_t
	- \alpha \frac{\hat{m}_t}{|g_t| + \varepsilon}.
	\]
	Since \(\hat{m}_t \approx g_t\) when variance is small, the expression simplifies to
	\[
	\theta_{t+1}
	\approx \theta_t - \alpha\, \text{sign}(g_t).
	\]
	Thus, with very small gradient variance, Adam behaves like a sign-based optimizer.
	
	\vspace{0.5em}
	\noindent\textbf{5(h).} Under the assumptions made here, show that the step size is scaling the gradient update to ensure it has unit variance.\\
	\noindent\textbf{Solution}\\
	With small gradient variance, \(\hat{m}_t \approx g_t\) and
	\(\hat{v}_t \approx g_t^2\). Then
	\[
	\frac{\hat{m}_t}{\sqrt{\hat{v}_t}}
	\approx \frac{g_t}{|g_t|}
	= \text{sign}(g_t),
	\]
	which has unit variance. Thus the step size rescales the gradient so that the
	update has variance one.
	
	\newpage
	
	\noindent\textbf{5(e)} Consider the optimal step-size derived for steepest-descent:
	\[
	f(x) = \frac{1}{2} x^{T} Q x - b^{T} x,
	\]
	\[
	\alpha_k = \frac{\nabla f_k^{T} \nabla f_k}{\nabla f_k^{T} Q \nabla f_k},
	\qquad
	x_{k+1} = x_k - \alpha_k \nabla f_k.
	\]
	
	Step schedulers reduce $\alpha_k$ when we approach convergence. Consider the
	simple case where we reduce the scheduler by a $c \alpha_k$ in epochs
	$e_1, e_2, e_3$. Assume models $Q_1, Q_2, Q_3$ for each one of the epochs.
	Derive expressions for $Q_2, Q_3$ based on $Q_1$. Sketch a 2-D example of what
	the models are assuming to be an ideal case for $c = 0.5$.\\
	\noindent\textbf{Solution}\\
	Reducing the step size by a factor $c$ gives the updated rule
	\[
	c\alpha_k
	= \frac{\nabla f_k^{T}\nabla f_k}{\nabla f_k^{T}Q_2\nabla f_k}.
	\]
	Matching this with the original formula implies
	\[
	Q_2 = \frac{1}{c}Q_1,
	\qquad
	Q_3 = \frac{1}{c^2}Q_1.
	\]
	For $c = 0.5$, this gives $Q_2 = 2Q_1$ and $Q_3 = 4Q_1$. The models therefore
	assume increasing curvature as the scheduler reduces the step size.
	
	\begin{figure}[h!]
		\centering
		\includegraphics[width=0.52\textwidth]{plots/fig_5.png}
		\caption{2-D sketch of $x^TQ_1x=1$, $x^TQ_2x=1$, and $x^TQ_3x=1$ for $c=0.5$.}
	\end{figure}
	
	\newpage
	
	\begin{lstlisting}[language=Python]
		import torch
		import torch.optim as optim
		import torch.nn as nn
		import matplotlib.pyplot as plt
		import numpy as np
		
		def get_loss(c_pred, c_true):
		return torch.mean((c_pred - c_true)**2)
		
		def predict(alpha0, alpha1, f):
		return alpha0 + alpha1 * f
		
		# ============================================================
		# Problem 3 – Automatic Differentiation
		# ============================================================
		
		x = torch.ones(6, requires_grad=True)
		f = (1/6) * torch.sum((x - 2)**2) + torch.sum(x)
		f.backward()
		print("Problem 3 gradient:", x.grad)
		
		# ============================================================
		# Problem 4(a) – Fahrenheit → Celsius using Autograd
		# ============================================================
		
		f_vals = 100 * torch.rand(10) + 32
		c_true = (f_vals - 32) / 1.8
		f_mean = f_vals.mean()
		f_std = f_vals.std()
		f_norm = (f_vals - f_mean) / f_std
		
		beta0 = torch.zeros(1, requires_grad=True)
		beta1 = torch.zeros(1, requires_grad=True)
		optimizer = optim.SGD([beta0, beta1], lr=0.1)
		
		max_iters = 2000
		loss = None
		
		for iteration in range(max_iters):
		optimizer.zero_grad()
		c_pred_norm = beta0 + beta1 * f_norm
		loss = get_loss(c_pred_norm, c_true)
		loss.backward()
		optimizer.step()
		
		with torch.no_grad():
		alpha1 = beta1 / f_std
		alpha0 = beta0 - beta1 * f_mean / f_std
		
		true_formula = "True:    C = -17.78 + 0.56·F"
		learned_formula = f"Learned: C = {alpha0.item():.2f} + {alpha1.item():.2f}·F"
		
		print("\nFahrenheit to Celsius")
		print("\t" + true_formula)
		print("\t" + learned_formula)
		print("Final loss:", loss.item())
		print("Iterations:", max_iters)
		
		# ============================================================
		# Problem 4(b) – Neural-Net Framework with SGD and Momentum
		# ============================================================
		
		class LinearNet(nn.Module):
		def __init__(self):
		super(LinearNet, self).__init__()
		self.linear = nn.Linear(1, 1)
		
		def forward(self, x):
		return self.linear(x)
		
		f_input_norm = f_norm.unsqueeze(1)
		c_target = c_true.unsqueeze(1)
		criterion = nn.MSELoss()
		
		net_sgd = LinearNet()
		optimizer_sgd = optim.SGD(net_sgd.parameters(), lr=0.1, momentum=0.0)
		num_epochs = 2000
		
		for epoch in range(num_epochs):
		optimizer_sgd.zero_grad()
		c_pred = net_sgd(f_input_norm)
		loss_sgd = criterion(c_pred, c_target)
		loss_sgd.backward()
		optimizer_sgd.step()
		
		with torch.no_grad():
		w_sgd = net_sgd.linear.weight.item()
		b_sgd = net_sgd.linear.bias.item()
		alpha1_sgd = w_sgd / f_std
		alpha0_sgd = b_sgd - w_sgd * f_mean / f_std
		
		net_mom = LinearNet()
		optimizer_mom = optim.SGD(net_mom.parameters(), lr=0.1, momentum=0.9)
		
		for epoch in range(num_epochs):
		optimizer_mom.zero_grad()
		c_pred_m = net_mom(f_input_norm)
		loss_mom = criterion(c_pred_m, c_target)
		loss_mom.backward()
		optimizer_mom.step()
		
		with torch.no_grad():
		w_mom = net_mom.linear.weight.item()
		b_mom = net_mom.linear.bias.item()
		alpha1_mom = w_mom / f_std
		alpha0_mom = b_mom - w_mom * f_mean / f_std
		
		print("\nProblem 4(b): Neural network with SGD")
		print("SGD (no momentum)")
		print(f"\tLearned: C = {alpha0_sgd.item():.2f} + {alpha1_sgd.item():.2f}·F")
		print(f"\tFinal loss: {loss_sgd.item()}")
		
		print("\nSGD (with momentum = 0.9)")
		print(f"\tLearned: C = {alpha0_mom.item():.2f} + {alpha1_mom.item():.2f}·F")
		print(f"\tFinal loss: {loss_mom.item()}")
		
		# ============================================================
		# Problem 5(e) – Q1, Q2, Q3 Sketch
		# ============================================================
		
		c = 0.5
		Q1 = np.array([[3, 1],[1, 2]])
		Q2 = Q1 / c
		Q3 = Q1 / (c**2)
		
		def plot_ellipse(Q, ax, label):
		theta = np.linspace(0, 2*np.pi, 400)
		unit_circle = np.vstack([np.cos(theta), np.sin(theta)])
		eigvals, eigvecs = np.linalg.eigh(Q)
		D_inv_sqrt = np.diag(1/np.sqrt(eigvals))
		transform = eigvecs @ D_inv_sqrt @ eigvecs.T
		ellipse = transform @ unit_circle
		ax.plot(ellipse[0], ellipse[1], label=label)
		
		fig, ax = plt.subplots(figsize=(6,6))
		plot_ellipse(Q1, ax, "Q1")
		plot_ellipse(Q2, ax, "Q2")
		plot_ellipse(Q3, ax, "Q3")
		ax.set_aspect('equal')
		ax.set_title("2D Sketch of Q1, Q2, Q3 for c = 0.5")
		ax.legend()
		plt.savefig("plots/fig_5.png")
		plt.show()
	\end{lstlisting}
	
\end{document}
